\subsection{Ideal functions on the bounded affine locus}
\label{subsec:ideal-functions}

Let $k$ be a trivially valued field, $A$ a commutative $k$-algebra and
$Y=\operatorname{Spec} A$. The ideal functions of
\cite[\S2.1]{FHX} are evaluated using generators in the local ring at a
semivaluation's center. The following affine statements impose precisely the
boundedness needed for independence of the generators. They do not construct
the projective Gauss extension or glue coherent ideal sheaves across charts.

\begin{definition}[Bounded points and ideal norms {\cite[\S2.1]{FHX}}]
\label{def:bounded-ideal-norm}
\lean{EnvelopeProperty.Berk.IsBounded, EnvelopeProperty.Berk.idealNorm,
EnvelopeProperty.Berk.idealLog, EnvelopeProperty.Berk.idealNorm_le_one,
EnvelopeProperty.Berk.idealLog_le_zero}
\leanok
\uses{def:berk,def:general-ideal-norm,spectrum:trivial-field-bridge}
A point $x\in Y^{\mathrm{an}}$ is bounded on $A$ if $|a|_x\le1$ for every
$a\in A$. For an ideal $I\subseteq A$ put
\[
 |I|_x=\sup_{a\in I}|a|_x\in[0,+\infty],\qquad
 \log|I|(x)=\log |I|_x\in[-\infty,+\infty].
\]
If $x$ is bounded then $|I|_x\le1$ and $\log|I|(x)\le0$.
\end{definition}

\begin{lemma}[Evaluation of generated ideals {\cite[\S2.1]{FHX}}]
\label{lem:ideal-norm-span}
\lean{EnvelopeProperty.Berk.idealNorm_span,
EnvelopeProperty.Berk.idealNorm_span_finset,
EnvelopeProperty.Berk.idealNorm_span_finset_eq_of_span_eq}
\leanok
\uses{def:bounded-ideal-norm,lem:general-ideal-norm-generators}
If $x$ is bounded and $S\subseteq A$, then
\[
 |(S)|_x=\sup_{s\in S}|s|_x.
\]
In particular, if $I=(f_1,\ldots,f_r)$ then
$|I|_x=\max_{1\le i\le r}|f_i|_x$, with the empty maximum equal to $0$.
Consequently this maximum is independent of the finite generating family.
\end{lemma}
\begin{proof}
Every $s\in S$ lies in $(S)$, giving the lower bound. Write
$M=\sup_{s\in S}|s|_x$. The set of $a\in A$ satisfying $|a|_x\le M$
contains $S$ and $0$, is closed under addition by the ultrametric inequality,
and is closed under multiplication by $b\in A$, since
$|ba|_x=|b|_x|a|_x\le|a|_x\le M$. It therefore contains $(S)$,
giving the upper bound. A supremum over a finite set is its maximum;
applying the equality to either generating family proves independence.
\end{proof}

\begin{lemma}[Ideal arithmetic {\cite[\S2.1]{FHX}}]
\label{lem:ideal-norm-arithmetic}
\lean{EnvelopeProperty.Berk.idealNorm_sup,
EnvelopeProperty.Berk.idealNorm_mul,
EnvelopeProperty.Berk.idealNorm_pow}
\leanok
\uses{lem:ideal-norm-span,lem:general-ideal-norm-arithmetic}
For a bounded point $x$ and ideals $I,J\subseteq A$,
\[
 |I+J|_x=\max\{|I|_x,|J|_x\},\qquad
 |IJ|_x=|I|_x|J|_x,\qquad |I^n|_x=|I|_x^n.
\]
The zero ideal has norm $0$ and the unit ideal has norm $1$.
\end{lemma}
\begin{proof}
View $x$ as a multiplicative seminorm on its coordinate ring. The generic
ideal arithmetic lemma gives the sum and product identities for this
seminorm, even without boundedness.
The norm of the zero ideal is $|0|_x=0$; the norm of the unit ideal is
at most $1$ by boundedness and at least $|1|_x=1$. Induction on $n$,
using the product equality and the unit ideal base case, proves the
power equality.
\end{proof}

\begin{lemma}[Logarithmic ideal formulas {\cite[\S2.1]{FHX}}]
\label{lem:ideal-log-arithmetic}
\lean{EnvelopeProperty.Berk.idealLog_sup,
EnvelopeProperty.Berk.idealLog_mul,
EnvelopeProperty.Berk.idealLog_pow,
EnvelopeProperty.Berk.idealLog_span_finset,
EnvelopeProperty.Berk.continuousOn_idealLog_of_span_eq}
\leanok
\uses{lem:ideal-norm-span,lem:ideal-norm-arithmetic,lem:general-ideal-log}
For a bounded point $x$,
\[
 \log|I+J|(x)=\max\{\log|I|(x),\log|J|(x)\},\qquad
 \log|IJ|(x)=\log|I|(x)+\log|J|(x),
\]
and $\log|I^n|(x)=n\log|I|(x)$.
For $I=(f_1,\ldots,f_r)$,
$\log|I|(x)=\max_i\log|f_i|(x)$, with the empty maximum equal
to $-\infty$. In particular the function determined by a finitely
generated ideal is continuous on the bounded affine locus.
\end{lemma}
\begin{proof}
The extended logarithm is order preserving and sends a product of
nonnegative extended real numbers to the sum of their logarithms.
Applying it to the preceding equalities gives the displayed formulas,
including the cases of a zero ideal norm. Its value at $0$ is
$-\infty$ and at $1$ is $0$, so the power formula includes $n=0$.
For finite generators it carries their maximum to the maximum of the
logarithms. Each function $x\mapsto\log|f_i|_x$ is continuous;
finite maxima of continuous functions are continuous, also for the
empty family, whose value is the constant $-\infty$.
\end{proof}

\begin{lemma}[Vanishing locus of an ideal function {\cite[\S2.1]{FHX}}]
\label{lem:ideal-function-vanishing}
\lean{EnvelopeProperty.Berk.idealNorm_eq_zero_iff,
EnvelopeProperty.Berk.idealLog_eq_bot_iff}
\leanok
\uses{def:bounded-ideal-norm,def:analytic-zero-locus,lem:analytic-affine-quotient,
lem:general-ideal-norm-generators}
For every affine analytic point $x$, without a boundedness assumption,
$|I|_x=0$ if and only if $I$ is contained in the support of $x$.
Equivalently $\log|I|(x)=-\infty$ exactly on the analytification of the
zero locus of $I$.
\end{lemma}
\begin{proof}
The supremum defining $|I|_x$ is zero if and only if every nonnegative
term $|a|_x$, $a\in I$, is zero. The support of the seminorm is its
zero ideal, so this condition is precisely $I\subseteq\operatorname{supp}(x)$.
The logarithm of a nonnegative extended real number equals $-\infty$
if and only if that number is zero.
\end{proof}

\begin{lemma}[Pullback of affine ideal functions {\cite[\S2.1]{FHX}}]
\label{lem:ideal-function-pullback}
\lean{EnvelopeProperty.Berk.IsBounded.comap,
EnvelopeProperty.Berk.idealNorm_map,
EnvelopeProperty.Berk.idealLog_map}
\leanok
\uses{lem:ideal-norm-span,lem:general-ideal-pullback}
Let $\varphi\colon A\to B$ be a $k$-algebra map and let $y$ be a
bounded point of $(\operatorname{Spec} B)^{\mathrm{an}}$. Its restriction
$x=\varphi^{\mathrm{an}}(y)$ is bounded on $A$. For the extended ideal
$IB=(\varphi(a):a\in I)$,
\[
 |IB|_y=|I|_x,\qquad \log|IB|(y)=\log|I|(x).
\]
\end{lemma}
\begin{proof}
For $a\in A$, $|a|_x=|\varphi(a)|_y\le1$, proving boundedness.
The extended ideal is generated by the images of all elements of $I$;
the generator formula gives
$|IB|_y=\sup_{a\in I}|\varphi(a)|_y=\sup_{a\in I}|a|_x$.
Taking logarithms gives the second equality.
\end{proof}
